\chapter{Building Commercial Applications with the Anthropic SDK and Claude Code SDK}
\label{app:sdk-commercial}

\begin{remark}[Learning Objectives]
After working through this appendix, the reader should be able to:
\begin{enumerate}
  \item Explain when the interactive Claude Code CLI is insufficient and a
        programmatic SDK approach is required for commercial deployments.
  \item Install the Anthropic Python or TypeScript SDK, authenticate with an
        API key, and make a basic completion call with the correct parameters.
  \item Define tools as JSON schemas, handle tool-call responses in a loop,
        and extract structured outputs using the SDK.
  \item Implement a self-contained agentic loop that manages conversation
        state, calls tools, and handles errors and retries without human input.
  \item Spawn and manage Claude Code processes headlessly via the Claude Code
        SDK, capturing structured output for downstream pipeline steps.
  \item Design a multi-agent workflow that fans work out in parallel, aggregates
        results, and passes state between orchestrator and subagents.
  \item Integrate token streaming into a web-facing API endpoint and handle
        partial response buffering correctly.
  \item Apply production-engineering practices---rate-limit handling, exponential
        backoff, cost budgeting, and structured logging---to an SDK-based service.
  \item Identify the security and compliance requirements specific to financial
        applications (API key hygiene, PII handling, audit trails, HITL controls).
  \item Trace through a complete automated financial report generator built on
        the Anthropic SDK, identifying where each design principle is applied.
\end{enumerate}
\end{remark}

% ---------------------------------------------------------------
\section{When to Go Beyond the Interactive CLI}
\label{app:sdk-commercial:motivation}

The interactive Claude Code CLI, described in Appendix~\ref{app:claude-code}, occupies
a well-defined niche: it is a force-multiplier for the individual developer who needs to
explore a codebase, draft a chapter, or run a one-off data analysis without leaving the
terminal. Its conversational interface, built-in tool integrations, and real-time display
of reasoning make it an exceptionally productive environment for tasks that benefit from
human steering. For that use case, it is close to optimal.

Commercial software products operate under fundamentally different constraints. A fintech
SaaS platform must serve thousands of concurrent users from a single code-base. A research
data pipeline must run nightly without human supervision, producing reports that feed
downstream risk systems. A real-time trading desk tool must respond in under a second to
incoming market data. None of these scenarios can be built on top of an interactive REPL.
They require programmatic, headless, scalable access to Claude's capabilities---access
provided by the Anthropic SDK and, for tasks that demand Claude Code's full agentic
toolkit, the Claude Code SDK \citep{anthropic2025api, anthropic2025ccsdk}.

This appendix introduces both interfaces. It begins with the limitations of the
interactive model for production systems, moves through the Anthropic SDK's core
primitives (messages, tool use, streaming), builds up to multi-agent orchestration
patterns, and concludes with the production-engineering and compliance concerns that any
commercial deployment in the financial industry must address. Throughout, code examples
are grounded in realistic financial use cases---earnings analysis, equity research
automation, and report generation---so that the architectural patterns are immediately
applicable.

\subsection{Limitations of the Interactive Model for Production Systems}
\label{app:sdk-commercial:motivation:limits}

The interactive CLI is, by design, a human-in-the-loop tool. Every invocation opens a
terminal session, renders output to a REPL, and expects a human to read and respond.
This design is a virtue for exploratory work and a fundamental constraint for
automation. Several concrete limitations emerge when one attempts to stretch the CLI
into a production role.

First, there is no headless operation mode. The CLI requires an attached terminal and,
implicitly, a human operator who can interpret its output and provide follow-up
instructions. Batch jobs, web servers, and microservices run without any attached
terminal; they must be invoked programmatically and must return results in a
machine-readable format. The CLI's primary output channel is a rich, human-readable
display that is not designed for downstream parsing.

Second, the CLI provides no mechanism for programmatic control over conversation state
or tool execution. A production system must be able to inject context into a prompt
dynamically (the current user's identity, account permissions, real-time market data),
route tool calls to specific implementations, intercept errors, and retry failed
operations---all without human input. The CLI abstracts these concerns away from the
user, but that abstraction becomes an obstacle when the system itself needs to be the
user.

Third, the interactive model cannot be embedded in multi-tenant web applications, API
servers, or CI/CD pipelines without resorting to fragile subprocess hacks that bypass
the CLI's output parsing, capture its stdout, and inject prompts through its stdin.
Such wrappers are brittle, hard to test, and carry no guarantees about compatibility
across CLI versions. The SDK, by contrast, is versioned, typed, and designed
specifically for programmatic use.

\begin{context}
A solo researcher running \texttt{/analyse earnings.pdf} to extract key metrics opens a
terminal, types a prompt, reads the response, and decides what to do next. The entire
workflow takes thirty seconds and requires one human. A fintech SaaS platform serving
ten thousand analysts cannot work this way. Each of those analysts submits requests
through a browser; the backend must call Claude, stream the response to the client,
log the exchange for compliance, and attribute the cost to the correct billing account,
all within a single HTTP request lifecycle. The interactive model is not the wrong tool
because it is bad; it is the wrong tool because it was designed for a different job.
The SDK exists precisely to fill this gap.
\end{context}

Fourth, the CLI offers no per-user isolation, cost attribution, or multi-tenancy. In a
commercial deployment, every API call must be associated with a user identity and an
organisational billing account. Rate limits must be enforced at the application layer.
Token spend must be tracked, capped, and reported. None of these concerns can be
addressed by the CLI without building the very abstraction layer that the SDK already
provides.

Finally, integration with CI/CD systems---automated test pipelines, pull request review
bots, documentation generators---requires that Claude be callable as a library function
returning a Python object or a TypeScript promise, not as a process whose output is
printed to a terminal. The SDK enables this integration pattern natively, while the CLI
requires subprocess plumbing that is fragile by comparison.

\subsection{Taxonomy of Programmatic Use Cases}
\label{app:sdk-commercial:motivation:taxonomy}

\begin{definition}[Programmatic LLM Integration]
\label{def:programmatic-llm}
A \emph{programmatic LLM integration} is any system in which calls to a large language
model are initiated by application code---not by a human interacting with a
conversational interface---and in which the model's responses are consumed by subsequent
code rather than being presented directly to a human for interpretation.
\end{definition}

Programmatic integrations span a wide range of architectural patterns. The following
taxonomy, while not exhaustive, covers the principal use cases encountered in financial
applications.

\begin{enumerate}
  \item \textbf{SaaS applications with embedded AI.} A web or mobile application
        integrates Claude to power features such as document summarisation, Q\&A over
        financial filings, or AI-assisted portfolio commentary. The application backend
        calls the SDK, streams results to the frontend, and bills the cost per user
        session.

  \item \textbf{Batch document processing pipelines.} A nightly job ingests hundreds of
        earnings transcripts, analyst reports, or regulatory filings, extracts structured
        data from each (revenue figures, forward guidance, risk factors), and loads the
        results into a data warehouse. Throughput and cost efficiency matter more than
        latency; the Haiku model tier is often appropriate here.

  \item \textbf{Automated research agents.} An agent loop retrieves financial data
        through tool calls, reasons over it across multiple turns, and produces a
        formatted research note. Unlike a batch pipeline, the agent's path through the
        task is not predetermined: it adapts based on what data it finds.

  \item \textbf{Real-time financial data analysis.} A market data feed triggers a Claude
        call whenever a watchlist stock moves more than two standard deviations in five
        minutes. The model's output---a one-paragraph situational summary---is pushed
        to a trader's dashboard within seconds. Streaming is essential here.

  \item \textbf{Internal tooling and APIs.} An internal REST API wraps Claude and
        exposes it to other teams as a service, with access controls, rate limits, and
        audit logging managed centrally. Individual teams consume the service without
        managing their own API keys or model parameters.

  \item \textbf{CI/CD code generation and review hooks.} A pre-merge hook calls Claude
        to review a pull request for compliance with internal coding standards, or to
        generate docstrings for new functions. The hook runs headlessly in a container,
        posts comments to the version control system, and fails the pipeline if critical
        issues are found.
\end{enumerate}

Each of these patterns places different demands on the SDK. Batch pipelines prioritise
throughput and cost; real-time dashboards prioritise latency and streaming; SaaS
applications require multi-tenancy and per-user cost attribution; CI/CD hooks require
determinism and structured output. The remainder of this appendix develops the
primitives and patterns that serve all of these needs.

% ---------------------------------------------------------------
\section{The Anthropic SDK}
\label{app:sdk-commercial:sdk}

The Anthropic SDK is the primary programmatic interface to all Claude models. It is
available as a Python package (\texttt{anthropic}) and a TypeScript/JavaScript package
(\texttt{@anthropic-ai/sdk}). Both packages wrap the same underlying HTTP API
\citep{anthropic2025api}, providing authentication, request serialisation, response
parsing, streaming support, automatic retries with exponential back-off, and typed
response objects that make it impossible to misread a field name at runtime.

The Python SDK is the natural choice for financial data science and research
applications, where the surrounding ecosystem---NumPy, pandas, scikit-learn, and
similar libraries---is predominantly Python. The TypeScript SDK is the right choice for
Node.js web servers, serverless functions, and any frontend-adjacent code that runs in
a JavaScript runtime. Throughout this appendix, primary examples are given in Python
with TypeScript equivalents noted where the idioms differ meaningfully.

The SDK abstracts away several concerns that would otherwise require substantial
boilerplate. Retry logic with jitter, connection pooling, chunked response parsing for
streaming events, and the translation of HTTP error codes into typed Python exceptions
are all handled automatically. A developer using the SDK can focus entirely on the
logic of their application rather than on the mechanics of the transport layer.

\subsection{Installation and Authentication}
\label{app:sdk-commercial:sdk:install}

Installation follows the standard package manager conventions for each ecosystem. In
Python, the SDK is distributed on PyPI:

\begin{minted}{bash}
pip install anthropic
\end{minted}

In Node.js and TypeScript environments, it is distributed on npm:

\begin{minted}{bash}
npm install @anthropic-ai/sdk
\end{minted}

Both packages require no native extensions; they are pure Python and pure
JavaScript respectively, which means they install without compilation on any
platform and work inside virtual environments, Docker containers, and serverless
runtimes without modification.

Authentication is managed through an API key issued from the Anthropic console. The
canonical pattern is to supply the key through the \texttt{ANTHROPIC\_API\_KEY}
environment variable, which both SDKs read automatically when no explicit key is
provided:

\begin{minted}{python}
import anthropic

# The client reads ANTHROPIC_API_KEY from the environment automatically.
# Do not pass the key as a literal string in source code.
client = anthropic.Anthropic()
\end{minted}

The injunction against hardcoding keys in source code is not merely a best practice; it
is a security requirement. API keys embedded in source files are trivially exposed
through version control history, container image layers, and log output. For local
development, store the key in a \texttt{.env} file that is listed in \texttt{.gitignore}
and load it with a library such as \texttt{python-dotenv}. In staging and production
environments, retrieve the key from a secrets manager---AWS Secrets Manager, HashiCorp
Vault, or GCP Secret Manager---at application startup and inject it into the process
environment. Never log the key, never pass it as a command-line argument, and never
include it in error messages.

\begin{minted}{python}
# Local development pattern using python-dotenv
from dotenv import load_dotenv
import os, anthropic

load_dotenv()  # reads .env into os.environ
client = anthropic.Anthropic(api_key=os.environ["ANTHROPIC_API_KEY"])
\end{minted}

The \texttt{Anthropic} constructor also accepts explicit \texttt{base\_url},
\texttt{max\_retries}, and \texttt{timeout} parameters, which are covered in
Section~\ref{app:sdk-commercial:production:rate-limits}.

\subsection{Making Your First API Call}
\label{app:sdk-commercial:sdk:first-call}

The core abstraction in the Anthropic API is the \emph{Messages} endpoint. A call to
\texttt{client.messages.create} sends a list of conversational turns to the model and
returns a \texttt{Message} response object. The minimal parameters are \texttt{model},
\texttt{max\_tokens}, and \texttt{messages}; the optional \texttt{system} parameter
provides a persistent instruction that frames the model's behaviour for the entire
conversation.

The following example asks Claude to distil the investment thesis for a hypothetical
company into a single sentence, a common operation in earnings analysis pipelines:

\begin{minted}{python}
import anthropic

client = anthropic.Anthropic()

message = client.messages.create(
    model="claude-opus-4-8",
    max_tokens=256,
    system="You are a financial analyst specialising in equity research.",
    messages=[
        {
            "role": "user",
            "content": (
                "In one sentence, summarise the investment thesis for a company "
                "reporting 20% revenue growth, expanding margins, and raising guidance."
            ),
        }
    ],
)

# The response text lives inside the first content block.
print(message.content[0].text)
print(
    f"Tokens used: {message.usage.input_tokens} in "
    f"/ {message.usage.output_tokens} out"
)
\end{minted}

The \texttt{Message} object returned by \texttt{create} has several fields that
production code regularly inspects. \texttt{message.content} is a list of content
blocks; for a simple text response there is exactly one block, an instance of
\texttt{TextBlock}, whose \texttt{.text} attribute holds the generated string.
\texttt{message.stop\_reason} reports why generation halted: \texttt{"end\_turn"}
means the model chose to stop; \texttt{"max\_tokens"} means the \texttt{max\_tokens}
limit was reached (a signal to increase the limit or split the request);
\texttt{"tool\_use"} means the model emitted a tool call requiring application-side
execution. \texttt{message.usage} is a \texttt{Usage} object with
\texttt{input\_tokens} and \texttt{output\_tokens} fields, both of which feed
directly into cost accounting.

\begin{example}[Typical token expenditure for financial summaries]
\label{ex:token-cost-summary}
A system prompt of 30 tokens plus a user message of 50 tokens totals 80 input tokens.
At the \texttt{claude-opus-4-8} rate of \$15 per million input tokens, this costs
\$0.0000012 per call. One million such calls---a plausible daily volume for a
high-traffic fintech SaaS---cost \$1.20 in input tokens plus output token costs
depending on response length. At 50 output tokens per call and \$75 per million output
tokens, the output cost is \$3.75 per million calls. Total cost per million calls is
approximately \$4.95, illustrating why model selection matters at scale.
\end{example}

\subsection{Model Selection and Parameters}
\label{app:sdk-commercial:sdk:params}

The Anthropic model family is organised into three capability tiers, each optimised for
a different point on the speed--cost--quality trade-off curve. Table~\ref{tab:model-tiers}
summarises the key characteristics.

\begin{table}[H]
  \centering
  \caption{Anthropic model tiers and indicative pricing (as of 2026).}
  \label{tab:model-tiers}
  \begin{tabular}{llrrp{4.5cm}}
    \toprule
    Model ID & Tier & Input (\$/M tok) & Output (\$/M tok) & Best suited for \\
    \midrule
    \texttt{claude-haiku-4-5-20251001} & Haiku & 0.25 & 1.25
      & Classification, tagging, entity extraction, high-volume batch jobs \\
    \texttt{claude-sonnet-4-6}         & Sonnet & 3.00 & 15.00
      & Analysis, summarisation, code generation, streaming chat \\
    \texttt{claude-opus-4-8}           & Opus & 15.00 & 75.00
      & Complex multi-step reasoning, long-context synthesis, research agents \\
    \bottomrule
  \end{tabular}
\end{table}

Beyond \texttt{model} and \texttt{max\_tokens}, the most consequential parameters are
\texttt{system}, \texttt{temperature}, and \texttt{stop\_sequences}. The system prompt
establishes the model's role, constraints, output format, and any domain knowledge that
should be assumed throughout the conversation; it is the principal lever for shaping
response behaviour and should be crafted with the same care as a software interface
contract. The \texttt{temperature} parameter, which ranges from zero to one with a
default of one, controls the stochasticity of token sampling: higher values produce
more varied and creative outputs, while lower values concentrate probability mass on the
most likely tokens. In financial applications, determinism is often more valuable than
creativity; setting \texttt{temperature=0} in production pipelines produces consistent
outputs that are easier to test and audit.

\begin{remark}
\label{rem:temperature-zero}
Setting \texttt{temperature=0} does not guarantee bit-for-bit identical outputs across
all runs, because the model may still encounter rounding differences at the hardware
level across different inference hosts. For tasks that require absolute reproducibility
(regression testing, audit trails), record the full response including the
\texttt{model} field and timestamp alongside the output. Do not rely on identical token
sequences; rely instead on semantic equivalence validated by a downstream parser or
schema check.
\end{remark}

The \texttt{stop\_sequences} parameter accepts a list of strings; generation halts as
soon as any of them appears in the output. This is useful for extracting the first
sentence of a response or for terminating generation once a delimiter such as
\texttt{"---"} is emitted. In multi-turn agent loops, stop sequences can signal the
model to pause and wait for tool results rather than continuing to hallucinate
downstream steps.

For tasks that are clearly within the Haiku tier's capabilities---single-label
classification, named-entity tagging, boolean yes/no judgements---the cost savings
relative to Opus are dramatic: at equal token counts, Haiku input tokens cost sixty
times less than Opus input tokens. A batch job that processes ten million earnings
transcript sentences costs \$2.50 in Haiku input tokens versus \$150 in Opus input
tokens. The discipline of matching model tier to task complexity is one of the highest-
leverage cost-reduction levers available to a production engineering team.

% ---------------------------------------------------------------
\section{Tool Use and Structured Outputs}
\label{app:sdk-commercial:tools}

The Messages API supports a mechanism by which the model can emit structured calls to
user-defined functions rather than generating free-form text. This capability,
sometimes called function calling or tool use, was studied empirically in the context
of general-purpose language models by \citet{schick2023toolformer}, who showed that
models can learn to insert API calls inline with generated text, and later extended to
the domain of API-calling by \citet{patil2023gorilla}, who demonstrated that models
can select correctly among thousands of real-world API signatures given natural
language descriptions. The Anthropic implementation follows the same conceptual pattern:
the developer declares a set of tools as JSON schemas, the model decides when to invoke
them based on the user's request, and the application executes the tool and returns the
result for the model to incorporate into its next response.

This mechanism is the foundation for building agents that interact with external systems.
A financial application might give Claude access to a data warehouse query function,
a market data feed, a portfolio management API, and a document retrieval system. The
model then acts as an intelligent router, selecting the right tool for each step of a
multi-turn reasoning process without the developer having to hard-code the routing logic.

\subsection{Defining Tools as JSON Schemas}
\label{app:sdk-commercial:tools:schemas}

Each tool is represented as a Python dictionary with three keys: \texttt{name} (a
machine-readable identifier), \texttt{description} (a natural-language explanation of
what the tool does and when to use it), and \texttt{input\_schema} (a JSON Schema
object that specifies the tool's parameters). The description is not merely
documentation; it is part of the model's context and directly influences the model's
decision about when and whether to call the tool. A vague description produces
unreliable tool selection; a precise description that explains both the tool's purpose
and its limitations produces robust behaviour.

The following example defines a tool for retrieving financial metrics from an internal
data warehouse, a common primitive in equity research applications:

\begin{minted}{python}
tools = [
    {
        "name": "get_financial_metric",
        "description": (
            "Retrieve a financial metric for a publicly traded company from the "
            "internal data warehouse. Use this whenever the user asks for specific "
            "financial figures such as revenue, EPS, P/E ratio, or market cap. "
            "Do not use this tool for qualitative questions or news-based queries; "
            "use it only for quantitative financial data that resides in the warehouse."
        ),
        "input_schema": {
            "type": "object",
            "properties": {
                "ticker": {
                    "type": "string",
                    "description": "Stock ticker symbol, e.g. AAPL, MSFT, NVDA.",
                },
                "metric": {
                    "type": "string",
                    "enum": [
                        "revenue", "eps", "pe_ratio",
                        "market_cap", "ebitda", "gross_margin",
                    ],
                    "description": "The financial metric to retrieve.",
                },
                "period": {
                    "type": "string",
                    "description": (
                        "Reporting period in the format Q1-2025 for quarterly data "
                        "or FY2024 for annual data. Omit for the most recent period."
                    ),
                },
            },
            "required": ["ticker", "metric"],
        },
    }
]
\end{minted}

Several design principles deserve emphasis. The \texttt{enum} constraint on
\texttt{metric} is deliberate: by restricting the model to a known set of metric names,
the application eliminates the risk of the model inventing a metric name that does not
exist in the warehouse schema, which would cause a silent lookup failure. Where
possible, enumerating valid values in the schema is preferable to accepting free-form
strings and validating them at runtime. The \texttt{required} list ensures the model
always provides the minimum necessary inputs; making \texttt{period} optional allows
the model to omit it and let the warehouse return the most recent value, which is the
natural behaviour for a user who asks ``what is Apple's EPS?'' without specifying a
period.

\subsection{Handling Tool Calls in the Response Loop}
\label{app:sdk-commercial:tools:loop}

When the model decides to call a tool, \texttt{stop\_reason} is set to
\texttt{"tool\_use"} and the \texttt{content} list contains one or more
\texttt{ToolUseBlock} objects in addition to any text blocks that precede the tool call.
Each \texttt{ToolUseBlock} has an \texttt{id} (a unique identifier for this specific
invocation), a \texttt{name} (matching one of the declared tool names), and an
\texttt{input} dict containing the arguments the model has chosen to pass.

The application must execute each tool call, collect the results, and return them to
the model as a new user-role message containing \texttt{tool\_result} content blocks.
Each result block references the tool call by its \texttt{id}, which allows the model
to correlate results with calls even when multiple tools are invoked in a single turn.

\begin{minted}{python}
import json
import anthropic

client = anthropic.Anthropic()

def get_financial_metric(ticker: str, metric: str, period: str | None = None) -> dict:
    """Stub: replace with actual warehouse query."""
    stub_data = {
        ("AAPL", "revenue"): {"value": 391_035_000_000, "currency": "USD",
                               "period": "FY2024"},
        ("AAPL", "eps"):     {"value": 6.43, "currency": "USD", "period": "FY2024"},
        ("MSFT", "revenue"): {"value": 245_122_000_000, "currency": "USD",
                               "period": "FY2024"},
    }
    key = (ticker.upper(), metric)
    return stub_data.get(key, {"error": f"No data for {ticker}/{metric}"})


def execute_tool(name: str, inputs: dict) -> dict:
    if name == "get_financial_metric":
        return get_financial_metric(**inputs)
    raise ValueError(f"Unknown tool: {name}")


def run_tool_loop(user_message: str) -> str:
    messages = [{"role": "user", "content": user_message}]

    while True:
        response = client.messages.create(
            model="claude-sonnet-4-6",
            max_tokens=1024,
            system="You are a financial analyst with access to a data warehouse.",
            tools=tools,
            messages=messages,
        )

        if response.stop_reason == "end_turn":
            # Return the final text response.
            return next(
                block.text
                for block in response.content
                if block.type == "text"
            )

        if response.stop_reason == "tool_use":
            # Append the assistant turn (which contains the tool call blocks).
            messages.append({"role": "assistant", "content": response.content})

            # Execute each tool call and collect results.
            tool_results = []
            for block in response.content:
                if block.type == "tool_use":
                    result = execute_tool(block.name, block.input)
                    tool_results.append({
                        "type": "tool_result",
                        "tool_use_id": block.id,
                        "content": json.dumps(result),
                    })

            # Return the results to the model as a user-role message.
            messages.append({"role": "user", "content": tool_results})
            continue

        # Unexpected stop reason; surface it.
        raise RuntimeError(f"Unexpected stop_reason: {response.stop_reason}")
\end{minted}

The loop above is deliberately simple: it alternates between assistant turns
(containing tool calls) and user turns (containing tool results) until the model
emits an \texttt{end\_turn}. In practice, the model may call multiple tools in a
single turn, and those calls may be interdependent. The application should execute
all calls in a given turn before returning results; if calls are independent, they
can be executed in parallel with \texttt{asyncio.gather} for lower latency.

\subsection{Structured Output with JSON Mode}
\label{app:sdk-commercial:tools:json}

A powerful pattern for extracting structured data from unstructured financial text is
to define a tool that encodes the desired output schema and then force the model to call
it using the \texttt{tool\_choice} parameter:

\begin{minted}{python}
import json, anthropic

client = anthropic.Anthropic()

extraction_tool = {
    "name": "extract_earnings_metrics",
    "description": "Extract key financial metrics from an earnings release or transcript.",
    "input_schema": {
        "type": "object",
        "properties": {
            "revenue_usd_millions":  {"type": "number"},
            "revenue_growth_yoy_pct": {"type": "number"},
            "eps_diluted":           {"type": "number"},
            "guidance_raised":       {"type": "boolean"},
            "key_risks":             {"type": "array",
                                     "items": {"type": "string"}},
        },
        "required": [
            "revenue_usd_millions",
            "revenue_growth_yoy_pct",
            "eps_diluted",
            "guidance_raised",
        ],
    },
}

def extract_from_earnings_text(text: str) -> dict:
    response = client.messages.create(
        model="claude-sonnet-4-6",
        max_tokens=512,
        system=(
            "You are a financial data extraction specialist. "
            "Extract metrics exactly as stated; do not infer or estimate values "
            "that are not explicitly present in the text."
        ),
        tools=[extraction_tool],
        tool_choice={"type": "tool", "name": "extract_earnings_metrics"},
        messages=[{"role": "user", "content": text}],
    )
    # With forced tool use, the first content block is always the ToolUseBlock.
    tool_block = response.content[0]
    return tool_block.input
\end{minted}

\begin{deepdive}
Forcing a specific tool call with \texttt{tool\_choice=\{"type": "tool", "name": ...\}}
is structurally more reliable than instructing the model to ``respond in JSON'' through
the system prompt. When JSON output is requested via prompting, the model may prefix the
JSON with a sentence of explanation, wrap it in a Markdown code fence, or omit
required fields when it has low confidence. Each of these behaviours requires custom
parsing logic that is brittle across model versions.

With forced tool use, the \texttt{input} field of the returned \texttt{ToolUseBlock}
is parsed and validated against the declared JSON Schema before it ever reaches
application code. The SDK raises a structured error if the model's output does not
conform to the schema, making validation failures explicit rather than silent. This
approach also eliminates the latency cost of parsing: there is no need to strip
Markdown fences, handle BOM characters, or detect whether the response is valid JSON
before attempting to parse it. The result is a typed Python dictionary that matches the
schema exactly, suitable for direct insertion into a database row or a pandas
DataFrame.
\end{deepdive}

% ---------------------------------------------------------------
\section{Building an Agent Loop}
\label{app:sdk-commercial:agent-loop}

The ReAct framework, introduced by \citet{yao2022react}, established the conceptual
architecture that underlies modern LLM agents: a loop in which the model alternates
between a \emph{reason} step---generating a chain-of-thought trace that analyses the
current state---and an \emph{act} step---emitting a tool call or a final answer.
In the Anthropic SDK, this loop maps naturally onto a while-loop over successive calls
to \texttt{client.messages.create}, with each iteration either appending a tool result
and continuing or returning the model's final text response.

The agent loop is the unit of autonomous action in agentic systems. Unlike a single
\texttt{messages.create} call, which completes a fixed exchange, an agent loop can
execute an arbitrary number of tool calls, accumulate evidence, revise hypotheses, and
eventually synthesise a conclusion---all without human intervention. This autonomy is
what makes the pattern useful for financial research tasks where the number and
sequence of data retrievals cannot be specified in advance.

\subsection{The Agentic Loop Pattern}
\label{app:sdk-commercial:agent-loop:pattern}

The canonical agent loop implementation wraps the tool-call logic from
Section~\ref{app:sdk-commercial:tools:loop} in a bounded iteration with a safety
guard against runaway execution. The \texttt{max\_turns} parameter is not optional in
a production system: without it, a model that repeatedly calls the wrong tool or
misinterprets tool results could exhaust a session's token budget before the
application notices.

\begin{minted}{python}
import json
import anthropic

client = anthropic.Anthropic()

def run_agent(
    user_prompt: str,
    tools: list,
    system: str = "You are a financial research agent with access to a data warehouse.",
    model: str = "claude-opus-4-8",
    max_turns: int = 10,
) -> str:
    """
    Run an agentic loop until the model emits end_turn or max_turns is reached.
    Returns the final text response.
    """
    messages = [{"role": "user", "content": user_prompt}]

    for turn in range(max_turns):
        response = client.messages.create(
            model=model,
            max_tokens=2048,
            system=system,
            tools=tools,
            messages=messages,
        )

        if response.stop_reason == "end_turn":
            for block in response.content:
                if block.type == "text":
                    return block.text
            return ""  # end_turn with no text block is unusual but safe to handle

        if response.stop_reason == "tool_use":
            messages.append({"role": "assistant", "content": response.content})
            tool_results = []
            for block in response.content:
                if block.type == "tool_use":
                    try:
                        result = execute_tool(block.name, block.input)
                    except Exception as exc:
                        result = {"error": str(exc)}
                    tool_results.append({
                        "type": "tool_result",
                        "tool_use_id": block.id,
                        "content": json.dumps(result),
                    })
            messages.append({"role": "user", "content": tool_results})
            continue

        raise RuntimeError(f"Unexpected stop_reason: {response.stop_reason}")

    raise RuntimeError(
        f"Agent exceeded max_turns={max_turns}. "
        "Increase the limit or check for a tool-call loop."
    )
\end{minted}

\begin{example}[Research agent for equity analysis]
\label{ex:research-agent}
A user asks: ``Compare Apple's and Microsoft's revenue growth over the last two fiscal
years and recommend which has the stronger top-line momentum.'' A single-turn call
cannot answer this: the model must first call \texttt{get\_financial\_metric} for
AAPL revenue in FY2024 and FY2023, then again for MSFT, accumulate four data points,
and only then synthesise a comparison. The agent loop handles this naturally: the model
plans the necessary retrievals, executes them across three to four turns, and then
generates a two-paragraph recommendation grounded in the retrieved data. The
developer's application code is the same regardless of how many tool calls the model
chooses to make.
\end{example}

\subsection{System Prompts and Conversation State}
\label{app:sdk-commercial:agent-loop:state}

The system prompt is the most stable component of the agent's context. It establishes
the model's role, its permissions (which tools it is authorised to use), its output
format requirements, and any domain constants that apply to all requests (the
organisation's investment mandate, the set of valid ticker symbols, the required tone
for client-facing prose). Because the system prompt is present in every turn, it is the
ideal candidate for prompt caching: the Anthropic API allows stable prefixes to be
marked with \texttt{cache\_control: \{type: ephemeral\}}, which reduces the input token
cost for subsequent calls that share the same prefix.

Dynamic context---the current date, the requesting user's identity, their portfolio
holdings, their risk tolerance profile---should be injected into the system prompt
using string formatting, but positioned after the stable portion so that caching is
not invalidated on every call. A common pattern is to construct the system prompt in
two parts: a stable preamble that describes the agent's role and tools, and a dynamic
suffix that provides session-specific context.

\begin{minted}{python}
from datetime import date

STABLE_SYSTEM = """You are a financial research agent for a long-only equity fund.

Your tools give you access to the fund's internal data warehouse, which contains
historical financial statements, consensus analyst estimates, and portfolio positions.
You should use these tools to ground your analysis in data rather than relying on
general knowledge, which may be stale.

When making a recommendation, always cite the specific data points that support it
and quantify the uncertainty where relevant. Do not fabricate data; if a requested
metric is unavailable, say so explicitly."""

def build_system_prompt(user_id: str, portfolio_id: str) -> str:
    dynamic_suffix = (
        f"\n\nSession context (do not reveal to the user):\n"
        f"- Current date: {date.today().isoformat()}\n"
        f"- User ID: {user_id}\n"
        f"- Portfolio ID: {portfolio_id}\n"
        f"- Authorised tickers: AAPL, MSFT, NVDA, GOOGL, AMZN, META\n"
    )
    return STABLE_SYSTEM + dynamic_suffix
\end{minted}

Conversation state is represented by the \texttt{messages} list, which grows with each
turn. For short agent loops---fewer than ten turns---holding the full conversation in
memory is fine. For long-running agents or sessions that span multiple HTTP requests
(in a stateless web server), the message list must be persisted externally, typically
in a Redis cache or a database table keyed by session ID. The SDK places no constraints
on how the list is stored; the application is free to serialise it to JSON and
deserialise it on the next request.

\subsection{Error Handling and Retries Within the Loop}
\label{app:sdk-commercial:agent-loop:errors}

Three categories of error arise in an agent loop, each requiring a different response
strategy. Understanding the distinctions prevents the common mistake of treating all
errors as equivalent and either swallowing them silently or halting the loop
prematurely.

The first category is API-level errors: network timeouts, connection resets, and rate
limit responses (HTTP 429). The Anthropic SDK retries these automatically with
exponential back-off up to \texttt{max\_retries} attempts (default four). Application
code should not attempt to retry these errors manually; doing so risks doubling the
retry budget and creating a thundering-herd effect under sustained rate pressure.
Instead, catch \texttt{anthropic.RateLimitError} only if the application needs to emit
a telemetry signal or apply additional back-pressure at the application layer.

The second category is tool execution errors: the tool function raises an exception
because the requested data does not exist, the warehouse is temporarily unavailable, or
the input parameters are invalid. These errors should not terminate the agent loop;
instead, they should be returned to the model as a structured error message in the
\texttt{tool\_result} content block. The model can then decide whether to try a
different approach, ask the user for clarification, or report the failure gracefully.

\begin{minted}{python}
import json
import anthropic

def safe_execute_tool(name: str, inputs: dict) -> str:
    """
    Execute a tool and return a JSON string.
    On failure, return a structured error rather than raising.
    """
    try:
        result = execute_tool(name, inputs)
        return json.dumps(result)
    except KeyError as exc:
        return json.dumps({"error": "not_found", "detail": str(exc)})
    except ConnectionError as exc:
        return json.dumps({"error": "warehouse_unavailable", "detail": str(exc)})
    except Exception as exc:
        return json.dumps({"error": "unexpected", "detail": str(exc)})
\end{minted}

The third category is model refusals or safety-filter activations, signalled by a
\texttt{stop\_reason} of \texttt{"stop\_sequence"} or by a response that contains no
text blocks and no tool calls. These are not errors in the SDK sense; they are
deliberate decisions by the model. The appropriate response is to log the
\texttt{stop\_reason} and any available metadata, surface a user-facing message that
the request could not be completed, and not retry the same prompt unchanged.

% ---------------------------------------------------------------
\section{The Claude Code SDK: Headless Execution}
\label{app:sdk-commercial:cc-sdk}

The Anthropic SDK and the Claude Code SDK address the same underlying need---
programmatic access to Claude---but at different levels of abstraction. The Anthropic
SDK exposes the raw Messages API: the developer calls the model, receives a response,
executes any tool calls, and manages the entire loop and tool ecosystem themselves.
The Claude Code SDK \citep{anthropic2025ccsdk}, by contrast, exposes a full Claude Code
agent headlessly: all of Claude Code's built-in tools---file reading, shell execution,
web search, MCP integrations---are pre-wired and available without any additional
implementation on the developer's part. The developer provides a prompt and a working
directory; the SDK handles everything else.

The choice between the two interfaces hinges on what the application needs to control.
If the task is calling a proprietary internal API, extracting structured data from a
custom schema, or implementing domain-specific tools, the Anthropic SDK is the right
choice because it gives full control over the tool ecosystem. If the task is agentic
coding, file manipulation, shell automation, or any other task that Claude Code already
handles well interactively, the Claude Code SDK is the right choice because it avoids
reimplementing that tooling.

\subsection{Architecture: The SDK Versus the Interactive CLI}
\label{app:sdk-commercial:cc-sdk:arch}

\begin{deepdive}
The interactive Claude Code CLI is a Node.js process that renders a rich terminal
interface using React and Ink. When the user types a prompt, the process sends it to
the Claude API, receives a stream of \texttt{SDKMessage} events, renders them to the
terminal in real time, and presents the result along with any tool outputs. The entire
interaction is mediated by a human reading the terminal.

The Claude Code SDK exposes the same underlying process in a non-interactive mode.
Internally, this is implemented by passing the \texttt{--print} flag (suppresses
the REPL) and \texttt{--output-format json} or \texttt{--output-format stream-json}
(routes output to stdout as structured JSON rather than a rendered terminal display).
The TypeScript SDK wraps this invocation, manages the child process lifecycle,
parses the JSON event stream, and exposes it to the caller as an async generator
that yields typed \texttt{SDKMessage} objects.

From the application's perspective, the difference is one of abstraction level. In
the CLI, the message stream is consumed by a renderer. In the SDK, the message stream
is consumed by the application. The underlying event types---\texttt{assistant},
\texttt{user}, \texttt{result}, \texttt{system}---are identical; only the consumer
changes. This means that any capability accessible in the interactive CLI is also
accessible through the SDK, including MCP server connections, custom tool
registrations, and configuration options passed through the agent's
\texttt{.claude/settings.json} file in the working directory.
\end{deepdive}

\subsection{Installation and Process Management}
\label{app:sdk-commercial:cc-sdk:install}

The Claude Code SDK is distributed as a Node.js package. Both the SDK and the Claude
Code CLI must be installed; the SDK delegates actual agent execution to the CLI binary.

\begin{minted}{bash}
# Install the CLI globally (required for the SDK to invoke it)
npm install -g @anthropic-ai/claude-code

# Install the SDK as a library dependency in a Node.js project
npm install @anthropic-ai/claude-code
\end{minted}

The SDK requires Node.js version 18 or later, which is the minimum version that
supports the \texttt{structuredClone} and \texttt{ReadableStream} APIs used internally.
Docker-based deployments should pin the Node.js version in the image definition to
avoid runtime surprises from upstream version changes.

The SDK manages the child process lifecycle automatically: it spawns the Claude Code
CLI as a child process when a query begins, streams its output through a pipe, and
terminates the process when the query completes or when the caller signals an abort via
an \texttt{AbortController}. Timeout management is the caller's responsibility; the
SDK does not enforce a default timeout. In production systems, always pass an
\texttt{AbortController} and call \texttt{abort()} after a configurable wall-clock
deadline to prevent runaway agent processes from consuming resources indefinitely.

For Python-based environments where installing a Node.js SDK is undesirable, the
subprocess approach described in Section~\ref{app:sdk-commercial:cc-sdk:spawn} is
the standard alternative.

\subsection{Spawning Tasks Programmatically}
\label{app:sdk-commercial:cc-sdk:spawn}

The TypeScript SDK's primary interface is the \texttt{query} async generator, which
accepts a prompt, an \texttt{AbortController}, and an options object specifying the
working directory, maximum turns, and model override. It yields a sequence of
\texttt{SDKMessage} objects; the final message of type \texttt{"result"} carries the
agent's terminal output as a string.

\begin{minted}{typescript}
import { query, type SDKMessage } from "@anthropic-ai/claude-code";

async function runCodeTask(
    prompt: string,
    cwd: string,
    maxTurns: number = 15,
    timeoutMs: number = 300_000,
): Promise<string> {
    const controller = new AbortController();
    const timeoutHandle = setTimeout(
        () => controller.abort(),
        timeoutMs,
    );

    let finalResult = "";

    try {
        for await (const message of query({
            prompt,
            abortController: controller,
            options: {
                maxTurns,
                cwd,
            },
        })) {
            if (message.type === "result") {
                finalResult = message.result;
            }
        }
    } finally {
        clearTimeout(timeoutHandle);
    }

    return finalResult;
}

// Example: generate a Python script that processes SEC filings.
const result = await runCodeTask(
    "Write a Python script that downloads the latest 10-K filing for AAPL "
    + "from the SEC EDGAR full-text search API and extracts the Risk Factors "
    + "section. Save the output to risk_factors.txt.",
    "/tmp/sec-analysis",
);
console.log(result);
\end{minted}

For Python environments, the same capability is accessible through a subprocess wrapper
that invokes the CLI directly and parses its JSON output:

\begin{minted}{python}
import subprocess
import json

def run_claude_code(
    prompt: str,
    cwd: str | None = None,
    max_turns: int = 10,
    timeout: int = 300,
) -> dict:
    """
    Run Claude Code headlessly and return the parsed JSON result.
    Raises RuntimeError if the process exits with a non-zero code.
    """
    cmd = [
        "claude",
        "--print",
        "--output-format", "json",
        "--max-turns", str(max_turns),
        prompt,
    ]
    proc = subprocess.run(
        cmd,
        cwd=cwd,
        capture_output=True,
        text=True,
        timeout=timeout,
    )
    if proc.returncode != 0:
        raise RuntimeError(
            f"Claude Code exited with code {proc.returncode}.\n"
            f"stderr: {proc.stderr}"
        )
    return json.loads(proc.stdout)
\end{minted}

The subprocess approach is straightforward but comes with one important limitation:
it cannot stream partial output to the caller, because \texttt{subprocess.run} waits
for the process to exit before returning. For long-running tasks---generating a large
codebase, analysing a multi-hundred-page PDF---this means the caller blocks for the
entire duration. For interactive applications this is unacceptable; for batch pipelines
it is usually tolerable. Section~\ref{app:sdk-commercial:cc-sdk:output} addresses the
streaming alternative.

\subsection{Capturing Structured Output}
\label{app:sdk-commercial:cc-sdk:output}

The \texttt{--output-format stream-json} flag instructs the CLI to emit each
\texttt{SDKMessage} as a newline-delimited JSON object as it is produced, rather than
accumulating them and emitting a single JSON blob at exit. This enables the Python
subprocess wrapper to read events incrementally using \texttt{subprocess.Popen} with
\texttt{stdout=PIPE}, processing each line as it arrives.

\begin{minted}{python}
import subprocess, json
from collections.abc import Generator

def stream_claude_code(
    prompt: str,
    cwd: str | None = None,
    max_turns: int = 10,
) -> Generator[dict, None, None]:
    """
    Yield SDKMessage dicts as they are emitted by Claude Code.
    The last message with type=="result" carries the final output.
    """
    cmd = [
        "claude",
        "--print",
        "--output-format", "stream-json",
        "--max-turns", str(max_turns),
        prompt,
    ]
    with subprocess.Popen(
        cmd,
        cwd=cwd,
        stdout=subprocess.PIPE,
        stderr=subprocess.PIPE,
        text=True,
    ) as proc:
        for line in proc.stdout:
            line = line.strip()
            if line:
                yield json.loads(line)
        proc.wait()
        if proc.returncode != 0:
            raise RuntimeError(proc.stderr.read())
\end{minted}

A common pattern in financial pipelines is to ask Claude Code to both implement a
computation and return a structured summary of its results. This is accomplished by
giving Claude an explicit instruction in the prompt: ``After running the analysis,
call the \texttt{report\_results} tool with the structured output.'' Combined with
a \texttt{.claude/settings.json} in the working directory that registers the tool
against a local HTTP endpoint, this creates a lightweight structured-output channel
that does not depend on parsing the agent's prose response.

% ---------------------------------------------------------------
\section{Multi-Agent Workflows}
\label{app:sdk-commercial:workflows}

Benchmarking studies on multi-agent systems, including \citet{liu2023agentbench}, have
demonstrated that collections of specialised agents consistently outperform single
general-purpose agents on complex, multi-domain tasks. The reasons are intuitive: a
single agent context window is finite, and a task that requires deep expertise in
multiple domains simultaneously will inevitably exceed what a single context can
accommodate. Decomposing the task across specialised agents, each focused on a narrow
sub-problem with a tailored system prompt and tool set, yields better results and
lower token costs because each agent's context is occupied only by information
relevant to its specific role.

In financial applications, the natural decomposition follows the structure of
professional investment teams: one agent for fundamental analysis, another for
technical analysis, a third for sentiment and news analysis, and an orchestrator that
synthesises the sub-analyses into a unified recommendation. This structure mirrors the
analyst team model used by sell-side research departments, which makes the resulting
system's outputs legible to the finance professionals who consume them.

\subsection{Orchestrator--Subagent Patterns}
\label{app:sdk-commercial:workflows:orchestrator}

\begin{definition}[Orchestrator--Subagent Pattern]
\label{def:orchestrator-subagent}
An \emph{orchestrator--subagent} system is a multi-agent architecture in which a
designated orchestrator agent receives the user's request, decomposes it into
subtasks, dispatches each subtask to a specialised subagent with a tailored prompt
and tool set, and aggregates the subagents' outputs into a final response. The
orchestrator is responsible for planning and synthesis; the subagents are responsible
for execution.
\end{definition}

The key design decision in an orchestrator--subagent system is the granularity of
decomposition. Decomposing too finely creates coordination overhead and multiplies the
number of API calls; decomposing too coarsely leaves individual agents with contexts
that are too large or roles that are too broad. For equity research, a three-subagent
decomposition---fundamentals, technicals, sentiment---is a natural starting point.

The orchestrator can be implemented either as a Claude agent (which plans dynamically
and may adjust the decomposition based on intermediate results) or as fixed application
code (which always fans out to the same set of subagents and then calls Claude for
synthesis). The fixed-code orchestrator is simpler to reason about, test, and debug;
the dynamic orchestrator is more flexible but harder to control. For production
financial applications where reliability and auditability are paramount, the fixed-code
orchestrator is usually the better choice.

\begin{example}[Equity research orchestrator]
\label{ex:orchestrator}
A user submits a request: ``Give me an investment recommendation for NVDA.'' The
application code dispatches three subagent calls in parallel (see
Section~\ref{app:sdk-commercial:workflows:parallel}): a fundamental subagent that
retrieves revenue, EPS, and margin data from the warehouse; a technical subagent
that retrieves price, volume, and moving average data; and a sentiment subagent that
retrieves recent news sentiment scores. Each subagent returns a 200-word analysis.
The orchestrator then calls the Opus model with all three analyses and asks for a
one-page investment recommendation. Total wall-clock time is dominated by the slowest
subagent call; with parallel execution, the three subagent calls contribute no more
latency than the slowest single call.
\end{example}

\subsection{Parallel Fan-Out with Async Processing}
\label{app:sdk-commercial:workflows:parallel}

Python's \texttt{asyncio} library and the \texttt{AsyncAnthropic} client make parallel
subagent execution straightforward. Each subagent call is an \texttt{async} function
that returns when the model responds; \texttt{asyncio.gather} runs all calls
concurrently and collects results when the last one completes.

\begin{minted}{python}
import asyncio
import anthropic

client = anthropic.AsyncAnthropic()

SUBAGENT_PROMPTS = {
    "fundamentals": (
        "Analyse the fundamental financial health of {ticker}: "
        "focus on revenue growth trend over the last three years, "
        "operating margin trajectory, and balance sheet leverage. "
        "Be specific and quantitative; cite the exact figures."
    ),
    "technicals": (
        "Analyse the technical price action of {ticker}: "
        "current trend (above or below 200-day MA), momentum (RSI), "
        "and key support/resistance levels. "
        "State a short-term directional bias."
    ),
    "sentiment": (
        "Analyse recent news and analyst sentiment for {ticker}: "
        "summarise the prevailing narrative, note any recent estimate revisions, "
        "and flag any material risks mentioned in the last 30 days."
    ),
}

async def run_subagent(ticker: str, aspect: str, prompt_template: str) -> dict:
    prompt = prompt_template.format(ticker=ticker)
    response = await client.messages.create(
        model="claude-haiku-4-5-20251001",
        max_tokens=512,
        system=(
            "You are a specialist financial analyst. "
            "Be concise, factual, and quantitative. "
            "Do not pad your response with disclaimers."
        ),
        messages=[{"role": "user", "content": prompt}],
    )
    return {
        "aspect": aspect,
        "analysis": response.content[0].text,
        "usage": {
            "input_tokens": response.usage.input_tokens,
            "output_tokens": response.usage.output_tokens,
        },
    }

async def parallel_equity_research(ticker: str) -> dict:
    # Fan out to all subagents simultaneously.
    tasks = [
        run_subagent(ticker, aspect, prompt)
        for aspect, prompt in SUBAGENT_PROMPTS.items()
    ]
    results = await asyncio.gather(*tasks)

    # Combine subagent reports for the orchestrator synthesis pass.
    combined = "\n\n".join(
        f"### {r['aspect'].title()} Analysis\n{r['analysis']}"
        for r in results
    )
    synthesis = await client.messages.create(
        model="claude-opus-4-8",
        max_tokens=1024,
        system=(
            "You are a senior portfolio manager at a long-only equity fund. "
            "Synthesise the subagent analyses into a concise investment recommendation "
            "with a BUY, HOLD, or SELL rating and a 12-month price target rationale."
        ),
        messages=[{
            "role": "user",
            "content": (
                f"Subagent analyses for {ticker}:\n\n{combined}\n\n"
                "Please synthesise these into a recommendation."
            ),
        }],
    )

    total_input = sum(r["usage"]["input_tokens"] for r in results)
    total_output = sum(r["usage"]["output_tokens"] for r in results)

    return {
        "ticker": ticker,
        "subagents": results,
        "recommendation": synthesis.content[0].text,
        "token_usage": {
            "subagent_input": total_input,
            "subagent_output": total_output,
            "orchestrator_input": synthesis.usage.input_tokens,
            "orchestrator_output": synthesis.usage.output_tokens,
        },
    }
\end{minted}

Using Haiku for the subagent calls and Opus only for the synthesis pass is a deliberate
cost optimisation. Each subagent receives a narrow, well-scoped prompt and produces a
short, factual paragraph; Haiku is fully capable of this task. The orchestrator
synthesis is the step that requires deep reasoning over multiple competing
considerations, which justifies the higher cost of Opus. This tier-differentiation
strategy---using the cheapest capable model for each sub-task---is a first-order
principle of cost-effective agentic system design.

\subsection{State Passing and Result Aggregation}
\label{app:sdk-commercial:workflows:state}

Inter-agent state passing requires careful attention to serialisation and scope. The
fundamental rule is that each subagent should receive only the information it needs to
complete its specific task---no more. Passing full conversation histories between agents
is wasteful (it fills the subagent's context with irrelevant turns), expensive (more
input tokens), and risks the ``telephone game'' failure mode, in which each layer of
summarisation introduces small distortions that compound into a significantly
misrepresented picture at the orchestrator.

\begin{remark}
\label{rem:telephone-game}
The ``telephone game'' error is one of the most insidious failure modes in multi-agent
systems. If subagent A summarises data before passing it to subagent B, and B
summarises that summary before passing it to the orchestrator, the orchestrator reasons
over a summary of a summary that may omit the very quantitative specifics needed for
a sound recommendation. The mitigation is to always pass raw data (the original numbers,
the original text) to each agent, and to perform summarisation only once, at the final
synthesis step.
\end{remark}

The preferred pattern is a typed dataclass that captures the subagent's output in a
structured form, with the raw data as a field alongside the prose analysis:

\begin{minted}{python}
from dataclasses import dataclass, field

@dataclass
class SubagentResult:
    aspect: str           # "fundamentals" | "technicals" | "sentiment"
    ticker: str
    analysis: str         # prose produced by the subagent
    raw_data: dict        # the warehouse data used (for audit purposes)
    input_tokens: int
    output_tokens: int

@dataclass
class OrchestratorResult:
    ticker: str
    recommendation: str   # "BUY" | "HOLD" | "SELL"
    rationale: str        # prose from the synthesis pass
    target_price: float | None
    subagent_results: list[SubagentResult] = field(default_factory=list)
    total_cost_usd: float = 0.0
\end{minted}

Storing the raw data in \texttt{SubagentResult.raw\_data} serves a dual purpose: it
enables the orchestrator to verify that the subagent's prose is consistent with the
underlying numbers (a lightweight hallucination check), and it provides a complete
audit record that can be replayed or independently verified after the fact.

% ---------------------------------------------------------------
\section{Streaming and Real-Time Applications}
\label{app:sdk-commercial:streaming}

Streaming is the mechanism by which the API delivers partial tokens to the client as
they are generated, rather than waiting for the entire response to be complete before
transmitting it. The practical effect is dramatic: a response that takes five seconds
to generate can begin appearing on the user's screen within two hundred milliseconds
of the request being submitted, because the first token arrives almost immediately and
subsequent tokens follow at the model's generation rate (typically twenty to eighty
tokens per second). For financial dashboards, analyst chat tools, and any other
application where a human is waiting for a response, streaming is the difference
between an application that feels responsive and one that feels broken.

The Anthropic SDK supports streaming through the \texttt{client.messages.stream()}
context manager (Python) and the \texttt{stream()} method on the messages object
(TypeScript). Both expose the token stream as an iterable of text fragments; the
application concatenates them in order to reconstruct the full response. Tool calls
are also streamed: the model emits incremental JSON fragments for each tool call's
\texttt{input} dict, which the SDK assembles into complete \texttt{ToolUseBlock}
objects before surfacing them to the application.

\subsection{Server-Sent Events and Token Streaming}
\label{app:sdk-commercial:streaming:sse}

The SDK's \texttt{stream()} context manager is the simplest entry point for streaming.
The \texttt{text\_stream} property yields text fragments as they arrive; the
\texttt{get\_final\_message()} method, available after the stream closes, returns the
complete \texttt{Message} object including usage statistics.

\begin{minted}{python}
import anthropic

client = anthropic.Anthropic()

def stream_analysis(ticker: str, question: str) -> None:
    """Stream a financial analysis to stdout, printing tokens as they arrive."""
    with client.messages.stream(
        model="claude-sonnet-4-6",
        max_tokens=1024,
        system="You are a financial analyst. Answer concisely and accurately.",
        messages=[{"role": "user", "content": f"{ticker}: {question}"}],
    ) as stream:
        for text in stream.text_stream:
            print(text, end="", flush=True)
    print()  # trailing newline after stream closes

    # Retrieve usage statistics after the stream is complete.
    final = stream.get_final_message()
    print(
        f"\n[Tokens: {final.usage.input_tokens} in "
        f"/ {final.usage.output_tokens} out]"
    )
\end{minted}

\begin{deepdive}
At the HTTP level, the Anthropic API uses Server-Sent Events (SSE), a lightweight
text-based protocol in which the server holds a single HTTP response open and
transmits newline-delimited \texttt{data:} lines as events occur. Each SSE line
carries a JSON payload whose \texttt{type} field identifies the event:
\texttt{content\_block\_start}, \texttt{content\_block\_delta},
\texttt{content\_block\_stop}, \texttt{message\_start}, \texttt{message\_delta},
and \texttt{message\_stop} are the principal types.

A \texttt{content\_block\_delta} event for text carries a payload of the form:
\begin{minted}{json}
{
  "type": "content_block_delta",
  "index": 0,
  "delta": {
    "type": "text_delta",
    "text": "Revenue grew 20"
  }
}
\end{minted}
The client concatenates all \texttt{delta.text} values for a given \texttt{index}
to reconstruct the complete text block. For tool calls, the delta type is
\texttt{input\_json\_delta} and the value is a JSON fragment that must be
concatenated and then parsed as a complete JSON object once the block closes.

Browsers can consume SSE natively using the \texttt{EventSource} API, which opens
a persistent HTTP connection and dispatches DOM events for each line received. This
makes SSE the natural choice for streaming LLM responses to a browser client:
the server proxies the Anthropic SSE stream, and the browser renders tokens in real
time using a simple event listener. The alternative---WebSockets---is more powerful
but requires a stateful server and a dedicated upgrade handshake; SSE is simpler and
sufficient for unidirectional token delivery.
\end{deepdive}

\subsection{Integrating Streaming into a Web API}
\label{app:sdk-commercial:streaming:web}

A FastAPI backend can proxy the Anthropic stream to a browser using the
\texttt{StreamingResponse} class with \texttt{media\_type="text/event-stream"}. The
generator function yields SSE-formatted lines as the Anthropic SDK delivers tokens;
the browser's \texttt{EventSource} or \texttt{fetch} with a \texttt{ReadableStream}
body reads them and appends each token to the displayed text.

\begin{minted}{python}
import json
import anthropic
from fastapi import FastAPI, Query
from fastapi.responses import StreamingResponse

app = FastAPI()
client = anthropic.Anthropic()

@app.get("/analyse/stream")
async def stream_endpoint(
    ticker: str = Query(..., description="Stock ticker symbol"),
    query: str = Query(..., description="Analysis question"),
):
    """
    Stream a financial analysis as Server-Sent Events.
    The client receives events of the form:
        data: {"token": "Revenue"}
        data: {"token": " grew"}
        ...
        data: [DONE]
    """
    async def token_generator():
        with client.messages.stream(
            model="claude-sonnet-4-6",
            max_tokens=2048,
            system=(
                "You are a financial analyst specialising in equity research. "
                "Answer the user's question about the named ticker concisely "
                "and with specific data where available."
            ),
            messages=[{
                "role": "user",
                "content": f"Ticker: {ticker}\nQuestion: {query}",
            }],
        ) as stream:
            for text in stream.text_stream:
                yield f"data: {json.dumps({'token': text})}\n\n"
        yield "data: [DONE]\n\n"

    return StreamingResponse(
        token_generator(),
        media_type="text/event-stream",
        headers={
            "Cache-Control": "no-cache",
            "X-Accel-Buffering": "no",  # disable nginx output buffering
        },
    )
\end{minted}

The \texttt{X-Accel-Buffering: no} header is important when the application sits
behind an nginx reverse proxy. By default, nginx buffers the upstream response until a
configurable threshold is reached, which defeats the purpose of streaming by
introducing a delay equal to the buffer-fill time. Setting this header instructs nginx
to forward each event to the browser immediately.

On the browser side, the stream can be consumed with the native \texttt{fetch} API
and a \texttt{ReadableStream} decoder:

\begin{minted}{typescript}
async function streamAnalysis(ticker: string, question: string): Promise<void> {
    const url = `/analyse/stream?ticker=${encodeURIComponent(ticker)}`
              + `&query=${encodeURIComponent(question)}`;
    const response = await fetch(url);
    const reader = response.body!.getReader();
    const decoder = new TextDecoder();
    let buffer = "";

    while (true) {
        const { done, value } = await reader.read();
        if (done) break;
        buffer += decoder.decode(value, { stream: true });
        const lines = buffer.split("\n\n");
        buffer = lines.pop() ?? "";  // keep the incomplete last chunk
        for (const line of lines) {
            if (line.startsWith("data: ")) {
                const payload = line.slice(6);
                if (payload === "[DONE]") return;
                const { token } = JSON.parse(payload);
                document.getElementById("output")!.textContent += token;
            }
        }
    }
}
\end{minted}

% ---------------------------------------------------------------
\section{Production Engineering}
\label{app:sdk-commercial:production}

The gap between a working prototype and a production-grade service is bridged by
systematic attention to a set of concerns that are invisible when running a notebook
interactively but become critical at scale: rate limits that truncate request bursts,
cost overruns that exhaust monthly budgets overnight, and silent failures that corrupt
downstream data without emitting any observable error signal. This section addresses
each concern in turn, providing concrete patterns that are directly applicable to
commercial financial deployments.

The Anthropic SDK provides several primitives that make production engineering easier:
configurable retry logic, typed exceptions for specific error categories, token
counting for pre-flight cost estimation, and streaming access to partial results.
The engineer's job is to compose these primitives into a robust service that handles
the full space of failure modes gracefully.

\subsection{Rate Limiting, Retries, and Back-Pressure}
\label{app:sdk-commercial:production:rate-limits}

Anthropic enforces rate limits at two levels: requests per minute (RPM) and tokens per
minute (TPM). Both limits vary by account tier; newly created accounts begin on a
conservative tier and can be upgraded by contacting Anthropic with a documented use
case. When a limit is exceeded, the API returns an HTTP 429 response, which the SDK
translates into an \texttt{anthropic.RateLimitError}.

By default, the SDK retries failed requests automatically, using exponential back-off
with jitter, up to \texttt{max\_retries} attempts (default four). For most applications
this is sufficient; the SDK's retry logic is well-tuned and the jitter prevents
correlated retry storms from a fleet of instances. The retry behaviour is configurable
at client construction time:

\begin{minted}{python}
import httpx
import anthropic

client = anthropic.Anthropic(
    max_retries=4,
    timeout=httpx.Timeout(60.0, connect=5.0),
)
\end{minted}

The \texttt{timeout} parameter accepts an \texttt{httpx.Timeout} object, which
separately controls the connection establishment timeout (how long to wait for a TCP
connection to the API server) and the read timeout (how long to wait for the next byte
of a response). For streaming responses where the model is generating a long output,
the read timeout should be set generously (sixty seconds or more) to avoid spurious
timeouts during slow generation; the connect timeout can remain short.

For applications that need application-level back-pressure---for example, a queue-based
batch processor that should slow down when the API is under pressure rather than
retrying indefinitely---the right pattern is to catch \texttt{RateLimitError} and
implement a backoff at the application layer:

\begin{minted}{python}
import time
import random
import anthropic

client = anthropic.Anthropic(max_retries=0)  # disable SDK retries; we manage them

def create_with_backoff(max_attempts: int = 5, **kwargs):
    """
    Wrapper around client.messages.create with manual exponential backoff.
    Suitable for queue consumers that want explicit control over retry pacing.
    """
    for attempt in range(max_attempts):
        try:
            return client.messages.create(**kwargs)
        except anthropic.RateLimitError:
            if attempt == max_attempts - 1:
                raise
            wait = (2 ** attempt) + random.uniform(0, 1)
            time.sleep(wait)
        except anthropic.APIStatusError as exc:
            if exc.status_code >= 500 and attempt < max_attempts - 1:
                wait = (2 ** attempt) + random.uniform(0, 1)
                time.sleep(wait)
            else:
                raise
\end{minted}

At the infrastructure level, a token-bucket rate limiter placed in front of the API
call layer (using Redis or a local in-memory implementation) provides a more
principled solution for high-throughput systems. The bucket is filled at the account's
TPM rate; each outgoing request drains the bucket by its estimated token count (from
\texttt{client.messages.count\_tokens}). Requests that would overdraw the bucket are
queued or shed according to the application's priority policy.

\subsection{Cost Management and Token Budgeting}
\label{app:sdk-commercial:production:cost}

Token-based billing creates a class of production risks that does not arise with
traditional software services: a single runaway request or a misconfigured
\texttt{max\_tokens} limit can generate an unexpectedly large bill. Three strategies
mitigate this risk at different levels of the stack.

The first is pre-flight token counting. The SDK's \texttt{client.messages.count\_tokens}
method accepts the same arguments as \texttt{client.messages.create} and returns the
number of input tokens that the request would consume, without actually sending it to
the model. This makes it possible to enforce a per-request input token budget before
any tokens are spent:

\begin{minted}{python}
import anthropic

client = anthropic.Anthropic()

class TokenBudget:
    """
    Tracks cumulative token spend across a session and enforces per-request
    and per-session limits.
    """
    def __init__(
        self,
        max_input_per_request: int = 50_000,
        max_input_per_session: int = 500_000,
        max_output_per_session: int = 100_000,
    ):
        self.max_input_per_request = max_input_per_request
        self.max_input_per_session = max_input_per_session
        self.max_output_per_session = max_output_per_session
        self.spent_input = 0
        self.spent_output = 0

    def pre_flight_check(self, model: str, messages: list, **kwargs) -> None:
        """Raise if this request would exceed per-request or per-session limits."""
        count = client.messages.count_tokens(
            model=model,
            messages=messages,
            **kwargs,
        )
        if count.input_tokens > self.max_input_per_request:
            raise RuntimeError(
                f"Request would use {count.input_tokens} input tokens, "
                f"exceeding per-request limit of {self.max_input_per_request}."
            )
        if self.spent_input + count.input_tokens > self.max_input_per_session:
            raise RuntimeError(
                f"Session input token budget exhausted: "
                f"{self.spent_input} used, "
                f"{count.input_tokens} requested, "
                f"{self.max_input_per_session} allowed."
            )

    def record(self, usage: anthropic.types.Usage) -> None:
        """Update cumulative spend after a successful API call."""
        self.spent_input  += usage.input_tokens
        self.spent_output += usage.output_tokens
        if self.spent_output > self.max_output_per_session:
            raise RuntimeError(
                f"Session output token budget exhausted: "
                f"{self.spent_output} output tokens used."
            )

    @property
    def estimated_cost_usd(self) -> float:
        """Estimated cost in USD at claude-sonnet-4-6 pricing."""
        return (
            self.spent_input  / 1_000_000 * 3.0
            + self.spent_output / 1_000_000 * 15.0
        )
\end{minted}

The second strategy is prompt caching. When a system prompt or a long document is
included in every call in a session, marking it with \texttt{cache\_control}
instructs Anthropic's infrastructure to cache the KV state for those tokens and charge
a discounted rate for cache hits on subsequent calls. The cache lasts five minutes
(``ephemeral'' tier) and is specific to the account. For an agent loop that appends a
large financial filing to every turn, caching the filing can reduce input token costs
by up to 90\% for all turns after the first.

\begin{minted}{python}
messages = [
    {
        "role": "user",
        "content": [
            {
                "type": "text",
                "text": long_earnings_transcript,
                "cache_control": {"type": "ephemeral"},  # cache this block
            },
            {
                "type": "text",
                "text": "What is the forward guidance for next quarter?",
            },
        ],
    }
]
\end{minted}

The third strategy is model tier selection, already discussed in
Section~\ref{app:sdk-commercial:sdk:params}. For a batch pipeline processing ten
thousand documents per day, switching from Sonnet to Haiku for the initial
classification pass (keeping Sonnet only for the documents that require deeper
analysis) can reduce the total bill by a factor of ten or more.

\subsection{Logging, Observability, and Tracing}
\label{app:sdk-commercial:production:observability}

Every LLM call in a production system should emit a structured log record that
captures sufficient information to replay, audit, and bill the call after the fact.
At minimum, this record should include: a unique request identifier, the model used,
the number of input and output tokens, the wall-clock latency, the \texttt{stop\_reason},
whether any tools were called (and which ones), and the estimated cost. For
multi-tenant applications, the user identifier and session identifier should also be
present.

\begin{minted}{python}
import time
import uuid
import logging
import dataclasses
from dataclasses import dataclass, field
from typing import Any

import anthropic

logger = logging.getLogger(__name__)

@dataclass
class LLMAuditRecord:
    request_id:     str
    user_id:        str
    session_id:     str
    model:          str
    input_tokens:   int
    output_tokens:  int
    latency_ms:     float
    stop_reason:    str
    tools_called:   list[str] = field(default_factory=list)
    estimated_cost_usd: float = 0.0
    error:          str | None = None

def _estimate_cost(model: str, input_tok: int, output_tok: int) -> float:
    pricing = {
        "claude-opus-4-8":           (15.0,  75.0),
        "claude-sonnet-4-6":         ( 3.0,  15.0),
        "claude-haiku-4-5-20251001": ( 0.25,  1.25),
    }
    in_rate, out_rate = pricing.get(model, (3.0, 15.0))
    return input_tok / 1_000_000 * in_rate + output_tok / 1_000_000 * out_rate

def instrumented_create(
    user_id: str,
    session_id: str,
    client: anthropic.Anthropic,
    **kwargs: Any,
) -> anthropic.types.Message:
    """
    Wrapper around client.messages.create that emits a structured audit record.
    """
    request_id = str(uuid.uuid4())
    start = time.perf_counter()
    error_msg: str | None = None
    response: anthropic.types.Message | None = None

    try:
        response = client.messages.create(**kwargs)
        return response
    except Exception as exc:
        error_msg = str(exc)
        raise
    finally:
        latency_ms = (time.perf_counter() - start) * 1000
        tools_called = []
        if response:
            tools_called = [
                b.name for b in response.content if b.type == "tool_use"
            ]
        record = LLMAuditRecord(
            request_id=request_id,
            user_id=user_id,
            session_id=session_id,
            model=kwargs.get("model", "unknown"),
            input_tokens=response.usage.input_tokens if response else 0,
            output_tokens=response.usage.output_tokens if response else 0,
            latency_ms=latency_ms,
            stop_reason=response.stop_reason if response else "error",
            tools_called=tools_called,
            estimated_cost_usd=_estimate_cost(
                kwargs.get("model", ""),
                response.usage.input_tokens if response else 0,
                response.usage.output_tokens if response else 0,
            ),
            error=error_msg,
        )
        logger.info("llm_call", extra={"audit": dataclasses.asdict(record)})
\end{minted}

For distributed systems in which a single user request fans out to multiple
concurrent Claude calls, distributed tracing with OpenTelemetry provides an
invaluable view of the call graph. Each Claude call should be instrumented as a child
span of the parent HTTP request span, with the request ID, model, and token counts
recorded as span attributes. The Anthropic SDK does not natively emit OpenTelemetry
spans, but wrapping each \texttt{instrumented\_create} call in a
\texttt{tracer.start\_as\_current\_span} context manager requires fewer than ten lines
of additional code and yields a complete distributed trace that correlates user
requests with individual model calls.

% ---------------------------------------------------------------
\section{Security and Compliance for Commercial Deployments}
\label{app:sdk-commercial:security}

Financial applications that incorporate LLMs face a regulatory environment
significantly more stringent than that faced by consumer-facing AI products. In
Europe, MiFID II requires that investment recommendations be made through a
documented and auditable process; AI-generated content does not exempt a firm from
this requirement. In the United States, FINRA Rule 3110 imposes supervisory obligations
on all business-related communications, including those generated or augmented by AI
systems. SEC guidance on AI tools in investment contexts emphasises the risk of
model-generated misinformation and the firm's duty to supervise AI outputs. GDPR and
CCPA impose constraints on the handling of personal data, which may include names,
account identifiers, and trading behaviour embedded in prompts sent to external APIs.

Compliance with this regulatory environment is not optional, and it cannot be
retrofitted cheaply after a system is built. The patterns described in this section
should be incorporated from the outset, not added as an afterthought.

\subsection{API Key Management and Secrets Hygiene}
\label{app:sdk-commercial:security:keys}

An API key is a credential that grants full access to an Anthropic account. Leaking a
key---through a committed \texttt{.env} file, a Docker image layer, a log line, or
an error message---is equivalent to handing an attacker a signed blank cheque against
the account's billing limit. The financial and reputational consequences of a key leak
can be severe; a well-funded adversary can exhaust a monthly token budget in minutes.

The canonical key-management pattern for production systems has four components. At
startup, the application retrieves the key from a secrets manager (AWS Secrets Manager,
HashiCorp Vault, Azure Key Vault, or GCP Secret Manager). The key is injected into the
process environment and passed to the \texttt{Anthropic} client constructor; it is
never stored in a file, a database, or an in-memory log buffer. The secrets manager
enforces access control (only the application's service account can retrieve the key)
and provides an audit trail of all retrieval events. The key is rotated on a
schedule---quarterly at minimum---and immediately on any suspected compromise, with the
old key revoked and a new key issued with no downtime gap.

\begin{minted}{python}
import os
import boto3  # AWS SDK; replace with equivalent for other cloud providers
import anthropic

def get_api_key_from_vault(secret_name: str, region: str = "eu-west-1") -> str:
    """Retrieve the Anthropic API key from AWS Secrets Manager."""
    client = boto3.client("secretsmanager", region_name=region)
    response = client.get_secret_value(SecretId=secret_name)
    # The secret value is a JSON string: {"ANTHROPIC_API_KEY": "sk-ant-..."}
    import json
    return json.loads(response["SecretString"])["ANTHROPIC_API_KEY"]

# Application startup sequence
api_key = get_api_key_from_vault("prod/anthropic-api-key")
llm_client = anthropic.Anthropic(api_key=api_key)
# api_key is now held only in the Anthropic client's internal state;
# do not log it, export it to other modules, or store it in instance variables.
\end{minted}

Local development environments should use a \texttt{.env} file that is explicitly
excluded from version control via \texttt{.gitignore}. A pre-commit hook that scans
staged files for strings matching the \texttt{sk-ant-} prefix provides a safety net
against accidental commits.

\subsection{Data Handling, PII, and Financial Regulations}
\label{app:sdk-commercial:security:pii}

Financial prompts regularly contain information that is legally sensitive: client
account numbers (PII), portfolio holdings (potentially subject to fund confidentiality
agreements), and in some contexts material non-public information (MNPI). Sending MNPI
to an external API raises regulatory concerns under insider trading regulations; even
where MNPI is not present, PII transmitted to a third party must be handled in
accordance with GDPR (for EU data subjects) or CCPA (for California residents).

Three mitigations are available, and they are best applied in combination. First,
establish a Data Processing Agreement (DPA) with Anthropic under an Enterprise plan
that explicitly excludes prompt content from use in model training. This does not
resolve MNPI concerns but does address the GDPR processor relationship requirement.
Second, apply a named-entity recognition (NER) step to prompts before they are sent,
replacing detected PII (names, account numbers, email addresses) with pseudonyms or
placeholder tokens:

\begin{minted}{python}
import re

def scrub_pii(text: str) -> str:
    """
    Remove or pseudonymise common PII patterns before sending to the API.
    This is illustrative; a production implementation should use a trained NER model.
    """
    # Replace account numbers (simplified pattern)
    text = re.sub(r'\b[A-Z]{2}\d{10,18}\b', '[ACCOUNT_NUMBER]', text)
    # Replace email addresses
    text = re.sub(
        r'\b[A-Za-z0-9._%+-]+@[A-Za-z0-9.-]+\.[A-Z|a-z]{2,}\b',
        '[EMAIL_ADDRESS]',
        text,
    )
    # Replace phone numbers (E.164 format and common variants)
    text = re.sub(r'\+?\d[\d\s\-().]{7,}\d', '[PHONE_NUMBER]', text)
    return text

def build_safe_prompt(raw_prompt: str) -> str:
    return scrub_pii(raw_prompt)
\end{minted}

Third, for the most sensitive workloads---hedge fund strategies, proprietary risk
models, client-specific recommendations---evaluate on-premises model deployment using
open-weight models (see Appendix~\ref{app:huggingface}) as a complement to or
replacement for the cloud API. On-premises deployment eliminates data egress entirely
at the cost of significantly higher infrastructure and operational overhead.

\subsection{Audit Trails and Human-in-the-Loop Controls}
\label{app:sdk-commercial:security:audit}

Regulatory obligations in financial services typically require that every AI-assisted
decision or recommendation be traceable after the fact: who requested it, when, what
the model was given, what it produced, and what action was taken. This requires an
audit log that is both comprehensive and tamper-resistant.

The audit record introduced in Section~\ref{app:sdk-commercial:production:observability}
provides the raw material. For regulatory purposes, this record should be written to an
append-only store: a database table with a write-once policy enforced at the schema
level (no \texttt{UPDATE} or \texttt{DELETE} permissions for the application service
account), or an immutable object store such as AWS S3 with Object Lock enabled. The
record should include the full prompt (after PII scrubbing), the full response, and the
hash of both (for integrity verification).

For high-stakes decisions---trade execution recommendations, credit approvals,
portfolio rebalancing proposals---a human-in-the-loop (HITL) gate must be inserted
between the model's output and any consequential action. The HITL gate is a review
queue: the model's recommendation is written to the queue with a \texttt{PENDING}
status, and the action is taken only after a human reviewer marks it \texttt{APPROVED}.
A FastAPI implementation of the HITL pattern:

\begin{minted}{python}
import uuid, datetime
from enum import Enum
from dataclasses import dataclass, field
from fastapi import FastAPI, HTTPException

app = FastAPI()

class ReviewStatus(str, Enum):
    PENDING   = "PENDING"
    APPROVED  = "APPROVED"
    REJECTED  = "REJECTED"

@dataclass
class PendingRecommendation:
    id:               str
    ticker:           str
    recommendation:   str   # "BUY" | "HOLD" | "SELL"
    rationale:        str
    target_price:     float | None
    created_at:       str
    status:           ReviewStatus = ReviewStatus.PENDING
    reviewed_by:      str | None = None
    reviewed_at:      str | None = None

# In-memory store for illustration; replace with a persistent database.
review_queue: dict[str, PendingRecommendation] = {}

@app.post("/recommendations/submit")
async def submit_recommendation(rec: dict) -> dict:
    """Submit an AI-generated recommendation for human review."""
    item = PendingRecommendation(
        id=str(uuid.uuid4()),
        ticker=rec["ticker"],
        recommendation=rec["recommendation"],
        rationale=rec["rationale"],
        target_price=rec.get("target_price"),
        created_at=datetime.datetime.utcnow().isoformat(),
    )
    review_queue[item.id] = item
    return {"id": item.id, "status": item.status}

@app.post("/recommendations/{rec_id}/approve")
async def approve_recommendation(rec_id: str, reviewer_id: str) -> dict:
    """Human reviewer approves a pending recommendation."""
    if rec_id not in review_queue:
        raise HTTPException(status_code=404, detail="Recommendation not found.")
    item = review_queue[rec_id]
    if item.status != ReviewStatus.PENDING:
        raise HTTPException(status_code=409, detail="Already reviewed.")
    item.status = ReviewStatus.APPROVED
    item.reviewed_by = reviewer_id
    item.reviewed_at = datetime.datetime.utcnow().isoformat()
    # Downstream: trigger trade execution, notify portfolio manager, etc.
    return {"id": rec_id, "status": item.status}
\end{minted}

The HITL gate does not eliminate the risk of acting on a flawed recommendation, but it
inserts a moment of deliberate human judgement that satisfies supervisory requirements
and provides a natural point at which the reviewer can flag anomalies for
investigation.

% ---------------------------------------------------------------
\section{Applied Example: An Automated Financial Report Generator}
\label{app:sdk-commercial:example}

This section assembles the primitives from the preceding sections into a complete,
end-to-end system. The system accepts a watchlist of stock tickers submitted by a user
via a REST endpoint, fans out to a set of analysis subagents, synthesises the results
into a formatted document, and returns a finished analyst report---all without human
intervention. It is a realistic, production-grade architecture; the components
described here are all used in commercial deployments of varying scale.

The system is intentionally modest in scope: it uses three tickers (AAPL, MSFT, NVDA),
a simplified data warehouse stub, and a text-only output format. Extending it to
support additional tickers, a real warehouse query layer, and PDF rendering via
\LaTeX\ or a reporting library is a matter of replacing the stubs with real
implementations; the orchestration logic is unchanged.

\subsection{Architecture Overview}
\label{app:sdk-commercial:example:arch}

The report generator is structured as a five-stage pipeline. In the first stage, the
user submits a watchlist of tickers through a POST endpoint. In the second stage, an
orchestrator function fans out to one analysis subagent per ticker, executing all
subagents in parallel using \texttt{asyncio.gather}. Each subagent retrieves financial
data through tool calls, reasons over the data, and returns a \texttt{ReportSection}
dataclass containing the analyst note, the recommendation, and the price target.

In the third stage, the orchestrator collects all sections and passes them to a
formatting agent (a single Claude call with the Sonnet model) that arranges the
sections into a coherent document structure with an executive summary and a ranking of
the tickers by recommendation strength. In the fourth stage, a validation pass
checks each section for internal consistency---data citations present, price targets
within plausible ranges---and flags any anomalies for the optional human review queue.
In the fifth stage, the validated document is returned to the caller as a structured
JSON object that can be rendered as HTML, a PDF, or a spreadsheet by the calling
application.

\begin{figure}[H]
  \centering
  \begin{tikzpicture}[
    node distance = 1.2cm and 2.2cm,
    box/.style = {rectangle, draw, rounded corners=3pt, minimum width=2.6cm,
                  minimum height=0.8cm, font=\small, align=center},
    arrow/.style = {->, >=stealth, thick},
  ]
    \node[box] (user)   {User\\Watchlist};
    \node[box, right=of user] (orch)  {Orchestrator};
    \node[box, above right=0.5cm and 1.6cm of orch] (sub1) {Subagent\\AAPL};
    \node[box, right=1.6cm of orch]                 (sub2) {Subagent\\MSFT};
    \node[box, below right=0.5cm and 1.6cm of orch] (sub3) {Subagent\\NVDA};
    \node[box, right=2.2cm of sub2]                 (fmt)  {Formatting\\Agent};
    \node[box, right=of fmt]                        (val)  {Validation\\Pass};
    \node[box, right=of val]                        (out)  {Report\\Output};

    \draw[arrow] (user) -- (orch);
    \draw[arrow] (orch) -- (sub1);
    \draw[arrow] (orch) -- (sub2);
    \draw[arrow] (orch) -- (sub3);
    \draw[arrow] (sub1) -- (fmt);
    \draw[arrow] (sub2) -- (fmt);
    \draw[arrow] (sub3) -- (fmt);
    \draw[arrow] (fmt)  -- (val);
    \draw[arrow] (val)  -- (out);
  \end{tikzpicture}
  \caption{Architecture of the automated financial report generator. The orchestrator
           fans out to one subagent per ticker in parallel; the formatting agent and
           validation pass run sequentially on the aggregated results.}
  \label{fig:report-gen-arch}
\end{figure}

\subsection{Implementation Walkthrough}
\label{app:sdk-commercial:example:impl}

The implementation begins with the data schema. Using typed dataclasses enforces
discipline about what each component produces and consumes, catches type mismatches at
development time, and makes the audit record self-documenting:

\begin{minted}{python}
from dataclasses import dataclass, field

@dataclass
class ReportSection:
    ticker:         str
    company_name:   str
    analyst_note:   str             # ~300-word prose analysis
    recommendation: str             # "BUY" | "HOLD" | "SELL"
    target_price:   float | None    # 12-month price target in USD
    raw_data:       dict = field(default_factory=dict)   # for audit
    input_tokens:   int = 0
    output_tokens:  int = 0

@dataclass
class FinancialReport:
    sections:            list[ReportSection]
    executive_summary:   str
    ranked_tickers:      list[str]  # best to least attractive
    total_input_tokens:  int = 0
    total_output_tokens: int = 0
    estimated_cost_usd:  float = 0.0
    validation_flags:    list[str] = field(default_factory=list)
\end{minted}

The analysis subagent is an async function that calls the Opus model with a finance-
specific system prompt and a structured user message containing the ticker's financial
data. The function parses the recommendation label from the response text using a
simple priority search and constructs a \texttt{ReportSection}:

\begin{minted}{python}
import asyncio, anthropic, json

client = anthropic.AsyncAnthropic()

async def analyse_ticker(ticker: str, data: dict) -> ReportSection:
    response = await client.messages.create(
        model="claude-opus-4-8",
        max_tokens=1024,
        system=(
            "You are a sell-side equity analyst at a bulge-bracket bank. "
            "Write a concise 300-word analyst note based strictly on the "
            "provided financial data. Structure your note as follows: "
            "one paragraph on revenue and growth trends, one paragraph on "
            "profitability and balance sheet, one paragraph on valuation and risks. "
            "End with a one-sentence investment recommendation containing the words "
            "BUY, HOLD, or SELL in capitals, followed by a 12-month price target "
            "in the format 'Target: $XXX'."
        ),
        messages=[{
            "role": "user",
            "content": (
                f"Ticker: {ticker}\n"
                f"Company: {data.get('name', ticker)}\n"
                f"Financial data:\n{json.dumps(data, indent=2)}"
            ),
        }],
    )

    text = response.content[0].text

    # Parse recommendation label.
    rec = "HOLD"
    for label in ["BUY", "SELL", "HOLD"]:
        if label in text.upper():
            rec = label
            break

    # Parse target price.
    import re
    target_match = re.search(r"Target:\s*\$(\d+(?:\.\d+)?)", text)
    target = float(target_match.group(1)) if target_match else None

    return ReportSection(
        ticker=ticker,
        company_name=data.get("name", ticker),
        analyst_note=text,
        recommendation=rec,
        target_price=target,
        raw_data=data,
        input_tokens=response.usage.input_tokens,
        output_tokens=response.usage.output_tokens,
    )


async def generate_report(tickers: list[str]) -> FinancialReport:
    # Fetch data for all tickers (stub; replace with warehouse query).
    data_map = {t: fetch_financial_data(t) for t in tickers}

    # Fan out to subagents.
    sections = list(await asyncio.gather(*[
        analyse_ticker(t, data_map[t]) for t in tickers
    ]))

    # Orchestrator synthesis: executive summary and ranking.
    combined = "\n\n".join(
        f"## {s.ticker}: {s.recommendation} | Target ${s.target_price}\n"
        f"{s.analyst_note}"
        for s in sections
    )
    synthesis_response = await client.messages.create(
        model="claude-sonnet-4-6",
        max_tokens=512,
        system=(
            "You are the head of equity research. Given individual analyst notes, "
            "write a two-paragraph executive summary of the overall outlook and "
            "rank the tickers from most to least attractive. "
            "Return a JSON object with keys 'executive_summary' (string) "
            "and 'ranked_tickers' (list of ticker strings)."
        ),
        tools=[{
            "name": "submit_report_metadata",
            "description": "Submit the executive summary and ticker ranking.",
            "input_schema": {
                "type": "object",
                "properties": {
                    "executive_summary": {"type": "string"},
                    "ranked_tickers":    {"type": "array",
                                         "items": {"type": "string"}},
                },
                "required": ["executive_summary", "ranked_tickers"],
            },
        }],
        tool_choice={"type": "tool", "name": "submit_report_metadata"},
        messages=[{"role": "user", "content": combined}],
    )

    meta = synthesis_response.content[0].input

    total_in  = sum(s.input_tokens  for s in sections)
    total_out = sum(s.output_tokens for s in sections)
    total_in  += synthesis_response.usage.input_tokens
    total_out += synthesis_response.usage.output_tokens
    cost = _estimate_cost("claude-opus-4-8", total_in, total_out)

    return FinancialReport(
        sections=sections,
        executive_summary=meta["executive_summary"],
        ranked_tickers=meta["ranked_tickers"],
        total_input_tokens=total_in,
        total_output_tokens=total_out,
        estimated_cost_usd=cost,
    )
\end{minted}

\subsection{Evaluation and Quality Gates}
\label{app:sdk-commercial:example:eval}

Before delivering the report to the user, a validation pass inspects each section for
common failure modes. Three checks are applied. The first verifies that each analyst
note contains at least one explicit reference to a data point from \texttt{raw\_data}:
a note that does not cite any figures may be hallucinating. The check is performed by
a Claude call that is given the note and the raw data and asked to confirm whether the
note is grounded.

The second check verifies that each price target is within a plausible range relative
to the current market price. A target that is more than 50\% above or below the current
price is flagged as a potential hallucination and routed to the human review queue
rather than included in the delivered report without comment. In practice, genuine
analyst targets occasionally exceed this range, but flagging them for review rather
than silently including them ensures that anomalous outputs receive scrutiny.

\begin{minted}{python}
def validate_section(section: ReportSection) -> list[str]:
    """
    Return a list of validation flags for the section.
    An empty list means the section passed all checks.
    """
    flags = []

    # Check 1: at least one specific data point cited.
    data_values = [str(v) for v in section.raw_data.values()
                   if isinstance(v, (int, float, str))]
    if not any(val[:6] in section.analyst_note for val in data_values):
        flags.append(
            f"{section.ticker}: analyst note may not cite specific data from warehouse."
        )

    # Check 2: price target plausibility.
    if section.target_price is not None:
        current_price = section.raw_data.get("current_price")
        if current_price and isinstance(current_price, (int, float)):
            ratio = section.target_price / current_price
            if ratio > 1.5 or ratio < 0.5:
                flags.append(
                    f"{section.ticker}: target price ${section.target_price:.2f} "
                    f"is outside ±50%% of current price ${current_price:.2f}. "
                    "Flagged for human review."
                )

    # Check 3: recommendation label present.
    if section.recommendation not in ("BUY", "HOLD", "SELL"):
        flags.append(
            f"{section.ticker}: recommendation label '{section.recommendation}' "
            "not recognised."
        )

    return flags

def validate_report(report: FinancialReport) -> FinancialReport:
    """Run validation on all sections and attach flags to the report."""
    all_flags = []
    for section in report.sections:
        all_flags.extend(validate_section(section))
    report.validation_flags = all_flags
    return report
\end{minted}

The third quality gate is a cost audit. After the report is generated, the orchestrator
logs the total token spend and estimated cost per report run. If the cost exceeds a
configurable threshold (for example, \$0.50 per report), an alert is raised so that
the engineering team can investigate whether a prompt has grown unexpectedly large or
whether a model has entered a high-token-consumption loop. Over time, cost-per-report
trends reveal whether system changes are moving costs in the intended direction and
whether the model-tier selection remains optimal as usage patterns evolve.

The complete system, assembled from the components above, processes a three-ticker
watchlist in approximately twelve to fifteen seconds of wall-clock time (dominated by
the parallel subagent calls, each of which takes four to six seconds for Opus at
current generation speeds). The estimated cost per run at three tickers is
approximately \$0.05 to \$0.10, depending on the length of the financial data provided
to each subagent. Scaling to twenty tickers increases wall-clock time by less than
50\% (because the subagents run in parallel) and cost linearly.
